\paragraph{Fixed-target posterior assessment.}
We assessed posterior computation on all \MTIFullSimulations{} simulation
samples and \MTIFullApplications{} application training partitions used below.
The selected \(q,\delta\), learning rate, and Gaussian root priors were held
fixed within each sample. Four marginal slice-within-Gibbs chains retained
4000 draws after 2000 burn-in cycles. Integrating out the scales gives the
same root posterior as the augmentation, with the conditional slice updates
described in the supplement \citep{Neal2003SliceSampling}.
All \MTIFullDatasets{} samples met the recorded diagnostic criteria for
ordered endpoints, width, and midpoint. The largest rank-normalized
\(\widehat R\) was \MTIFullRhat{}, and the smallest bulk and tail effective
sample sizes were approximately \MTIFullBulkESS{} and \MTIFullTailESS{},
respectively \citep{VehtariEtAl2021Rhat}.

For a point-interval comparison, let
\(I_{\mathrm G}=[\bar l,\bar u]\), where
\(\bar l=\E\{\min(a,b)\mid D_n,q,\delta\}\) and
\(\bar u=\E\{\max(a,b)\mid D_n,q,\delta\}\) under the generalized posterior.
Its width is the posterior mean of \(|b-a|\). We approximate these expectations
by the retained draws. This summary incorporates the spread of the root
posterior, while the finite ECM calculation estimates a mode. Empirical
limiting tilts can leave one outward root direction weakly constrained by
the loss, producing substantial width uncertainty under the diffuse prior.
Section~\ref{sec:evaluation} compares the resulting intervals on the same
validation samples; the supplement gives their full posterior and
repeated-sampling assessments.
